\chapter{Conclusion and Future Work}

\section{Summary of Contributions}

This thesis delivered \TimeLens{}, an AI- and AR-powered bilingual guide for the
\GEM{}, through two engineering contributions and a market study. The first is an
\textbf{on-device artifact detector} whose development isolated label quality---not
architecture or resolution---as the decisive factor, yielding a compact YOLOv8n model
that resolves every previously failing class. The second is a \textbf{grounded
bilingual RAG guide}, selected from a seven-model benchmark and optimized from over
30\,s to about 10\,s end-to-end latency, that keeps answers anchored to curated museum
data in English or Arabic. Both ship in a production Flutter application, and a market
analysis situates the system within the heritage-AR landscape.

\section{Critical Reflection}

The decisive gains in this project came not from model complexity but from
\textbf{clean data and grounding}: hand-verified labels made the on-device detector
trustworthy, and retrieval grounding kept the language model factual. The honest
residual risks are equally clear---validation frame leakage and training-to-deployment
domain shift on the detection side, fine-grained long-tail classes, and evolving
real-world capture conditions.

\section{Future Work}

\paragraph{Detection.} Split the dataset by source video rather than by frame to remove
frame leakage; capture and annotate a phone-camera \emph{deployment-distribution} test
set for honest field metrics; apply stronger augmentation for domain shift; and
re-label or expand the remaining long-tail classes.

\paragraph{Retrieval-augmented guide.} Expand the closed knowledge base; replace or
augment the keyword query classifier with a learned intent model; add conversational
memory across turns; and build index health-audit and maintenance tooling so the
knowledge base can be updated and validated routinely as the collection grows.

\paragraph{System.} Optionally enable backend re-verification (the \texttt{/predict}
path) for low-confidence detections; add explainable detection overlays; and conduct
primary user studies (including System Usability Scale evaluation) in live museum
conditions to quantify learning outcomes and long-horizon experience quality.
